% apendices/conjugacion-n4-n3.tex
\chapter{\ruby{活用表}{かつようひょう} N4--N3 — Tablas de Conjugación}

\section*{\ruby{動詞}{どうし}の\ruby{活用}{かつよう} — Conjugación Verbal}

\begin{tcolorbox}[vocabbox, title={G1 (Grupo 1 / 五段動詞) — \ruby{書}{か}く}]
\begin{tabularx}{\linewidth}{lXX}
\toprule
\textbf{Forma} & \textbf{書く (kaku)} & \textbf{Uso} \\ \midrule
\ruby{辞書形}{じしょけい} & \ruby{書}{か}く & forma de diccionario \\
ます形 & \ruby{書}{か}きます & educado \\
て形 & \ruby{書}{か}いて & secuencia / gerundio \\
た形 & \ruby{書}{か}いた & pasado casual \\
ない形 & \ruby{書}{か}かない & negación casual \\
ば形 & \ruby{書}{か}けば & condicional ば \\
させる形 & \ruby{書}{か}かせる & causativa \\
られる形 & \ruby{書}{か}かれる & pasiva \\
させられる & \ruby{書}{か}かせられる & causativa-pasiva \\
意向形 & \ruby{書}{か}こう & volición \\
命令形 & \ruby{書}{か}け & imperativo \\
\bottomrule
\end{tabularx}
\end{tcolorbox}

\begin{tcolorbox}[vocabbox, title={G2 (Grupo 2 / 一段動詞) — \ruby{食}{た}べる}]
\begin{tabularx}{\linewidth}{lXX}
\toprule
\textbf{Forma} & \textbf{食べる (taberu)} & \textbf{Uso} \\ \midrule
\ruby{辞書形}{じしょけい} & \ruby{食}{た}べる & forma de diccionario \\
ます形 & \ruby{食}{た}べます & educado \\
て形 & \ruby{食}{た}べて & secuencia / gerundio \\
た形 & \ruby{食}{た}べた & pasado casual \\
ない形 & \ruby{食}{た}べない & negación casual \\
ば形 & \ruby{食}{た}べれば & condicional ば \\
させる形 & \ruby{食}{た}べさせる & causativa \\
られる形 & \ruby{食}{た}べられる & pasiva / potencial \\
させられる & \ruby{食}{た}べさせられる & causativa-pasiva \\
意向形 & \ruby{食}{た}べよう & volición \\
命令形 & \ruby{食}{た}べろ & imperativo \\
\bottomrule
\end{tabularx}
\end{tcolorbox}

\begin{tcolorbox}[vocabbox, title={G3 (Irregulares) — する / くる}]
\begin{tabularx}{\linewidth}{lXXX}
\toprule
\textbf{Forma} & \textbf{する} & \textbf{くる} & \textbf{Uso} \\ \midrule
\ruby{辞書形}{じしょけい} & する & くる & diccionario \\
ます形 & します & きます & educado \\
て形 & して & きて & gerundio \\
た形 & した & きた & pasado \\
ない形 & しない & こない & negación \\
ば形 & すれば & くれば & condicional \\
させる & させる & こさせる & causativa \\
られる & される & こられる & pasiva/pot. \\
させられる & させられる & こさせられる & caus.-pasiva \\
意向形 & しよう & こよう & volición \\
\bottomrule
\end{tabularx}
\end{tcolorbox}

\section*{\ruby{形容詞}{けいようし}の\ruby{活用}{かつよう} — Conjugación de Adjetivos}

\begin{tcolorbox}[vocabbox, title={い形容詞 — \ruby{高}{たか}い / な形容詞 — \ruby{静}{しず}かだ}]
\begin{tabularx}{\linewidth}{lXX}
\toprule
\textbf{Forma} & \textbf{高い (alto)} & \textbf{静かだ (silencioso)} \\ \midrule
\ruby{辞書形}{じしょけい} & \ruby{高}{たか}い & \ruby{静}{しず}かだ \\
丁寧形 & \ruby{高}{たか}いです & \ruby{静}{しず}かです \\
過去形 & \ruby{高}{たか}かった & \ruby{静}{しず}かだった \\
否定形 & \ruby{高}{たか}くない & \ruby{静}{しず}かじゃない \\
て形 & \ruby{高}{たか}くて & \ruby{静}{しず}かで \\
ば形 & \ruby{高}{たか}ければ & \ruby{静}{しず}かであれば \\
たら形 & \ruby{高}{たか}かったら & \ruby{静}{しず}かだったら \\
副詞形 & \ruby{高}{たか}く & \ruby{静}{しず}かに \\
\bottomrule
\end{tabularx}
\end{tcolorbox}

\section*{\ruby{て形}{てけい}の\ruby{接続表現}{せつぞくひょうげん} — Expresiones con て形}

\begin{tcolorbox}[vocabbox, title={て + 補助動詞 — Verbos auxiliares después de て}]
\begin{tabularx}{\linewidth}{lXX}
\toprule
\textbf{Patrón} & \textbf{Significado} & \textbf{Casual} \\ \midrule
〜ている & estado continuo / resultante & 〜てる \\
〜てしまう & completar (con matiz) & 〜ちゃう / 〜じゃう \\
〜ておく & hacer de antemano & 〜とく \\
〜てある & resultado de acción deliberada & — \\
〜てみる & intentar / probar & — \\
〜てくる & cambio que se acerca & — \\
〜ていく & cambio que se aleja & — \\
〜てほしい & querer que otro haga & — \\
〜てもいい & tener permiso & — \\
〜てはいけない & estar prohibido & 〜ちゃだめ \\
\bottomrule
\end{tabularx}
\end{tcolorbox}
